% --- Archon physics formalization source begin ---
% archon:physics
% archon:chemistry
% archon:covers IChO2026Problems/problem_icho_2026_t3_a1.lean
% archon:source-report reports/icho_2026/problem_icho_2026_t3_a1.source.json
% archon:problem-id icho_2026_t3
% archon:part-id T3-A1

\chapter{Icho Chemistry problem icho\_2026\_t3\_a1}
\label{ch:IChO2026Problems_problem_icho_2026_t3_a1}

\paragraph{Problem source.}
\#\# Source
58th International Chemistry Olympiad, Tashkent, Uzbekistan, 2026.
Official-materials index: https://www.icho2026.uz/problems
English problem PDF: ../icho\_2026\_source/raw/theory\_problem.pdf
English solution/rubric PDF: ../icho\_2026\_source/raw/theory\_solution.pdf
Problem PDF page 25; printed local question page 1; solution starts on PDF page 22.
The page PNG is primary visual evidence for formulas, structures, tables, and figures.

\#\# T3. Into Reticular Chemistry
T3. Into Reticular Chemistry 7\% Covalent organic frameworks (COFs) are a class of polymers that form two- or three-dimensional structures through reactions between organic precursors, resulting in covalent bonds to afford porous, crystalline materials. Two examples of honeycomb type 2D-COFs formed by boronate ester and boroxine rings (repeat units indicated using dashed lines) are shown in the figure below. The black and red bold lines represent different building blocks in all COF topology figures.

\#\# Current subquestion 3.1
Determine the empirical formula of COF-1 and calculate the mass percentage (\%, to two decimal places) of carbon in COF-1.

\#\# Dataset metadata
IChO 2026; T3; raw rubric points=2; kind=theory; formalization\_ready=true.

\paragraph{Shared problem context.}
T3. Into Reticular Chemistry 7\% Covalent organic frameworks (COFs) are a class of polymers that form two- or three-dimensional structures through reactions between organic precursors, resulting in covalent bonds to afford porous, crystalline materials. Two examples of honeycomb type 2D-COFs formed by boronate ester and boroxine rings (repeat units indicated using dashed lines) are shown in the figure below. The black and red bold lines represent different building blocks in all COF topology figures.

\paragraph{Current subquestion.}
Determine the empirical formula of COF-1 and calculate the mass percentage (\%, to two decimal places) of carbon in COF-1.

\paragraph{Recorded answer/context.}
The empirical formula of COF-1 is C9H4BO2, its elemental composition is C – 12.01⋅9 12.01⋅9+4⋅1.008+10.81+32 ⋅100\% = 69.77\%; C: 69.77\% Empirical formula: C9H4BO2 1 pt for empirical formula. No penalty for [C9H4BO2]𝑛 1 pt for C mass fraction. no penalty if integer values are used for atomic masses 2.0 pt

\paragraph{Figure/image path.}
../icho\_2026\_source/image/T3\_page-1.png

\paragraph{Formalization target.}
The verified Lean formalization is in `IChO2026Problems/problem\_icho\_2026\_t3\_a1.lean`,
with its theorem and lemma proofs complete.
The Lean declarations must preserve every mathematical object, quantifier, hypothesis, side condition, approximation/error guarantee, and requested conclusion from the source statement.
The implementation reuses the pinned Mathlib, Physlib, and Chemistry libraries,
with local definitions only for source-specific objects.

\begin{theorem}[Icho Chemistry formalization target]
\lean{IChO2026Problems.Icho2026T3A1.cof1_empirical_formula_and_carbon_mass_percentage}
\leanok
\label{thm:physics:icho_2026_t3_a1:target}
The mapped declarations encode the stated calculation and its verified conclusion.
\end{theorem}
\begin{proof}

\leanok
The proof unfolds the source-specific definitions and establishes the mapped
conclusion from the encoded hypotheses.
\end{proof}
% --- Archon physics formalization source end ---
